\reportheading{Appendix F.\quad FDS inputs for two representative runs}
The propane example is the OP-6 coarse-grid benchmark at 106.2~kW. The propylene example is the OP-6 run at soot yield 0.05 that governs thermal clearance at the selected 1.12~m$^3$/s. Both used FDS 6.10.1, $27.5\,^{\circ}$C ambient air and an adiabatic default surface. Coordinates are in metres and measured from the model origin.

\tabletitle{Table F1. Input settings checked against the two decks and solver records.}
\begin{reporttable}{L{.22}L{.38}L{.40}}
\rowcolor{tablegrey}Setting & Propane benchmark & Propylene design\\
Mesh and radiation & $132\times28\times19$, radiative fraction 0.29 & $132\times28\times38$, radiative fraction 0.37\\
Combustion & \texttt{PROPANE}, nominal 106.2~kW & \texttt{PROPYLENE}, nominal 106.2~kW. Soot and CO yields are 0.05 and 0.01~kg/kg fuel.\\
End-wall openings & Compartment ends closed & \texttt{OPEN} at $x=0$ and 6.60, across $y=0$--0.90 and $z=0$--0.48\\
Smoke curtains & None & At $x=2.80$ and 3.80, across $y=0$--0.90 and $z=0.52$--0.76\\
Ceiling extraction & None & Left vent $x=2.90$--3.30, right vent $x=3.30$--3.70. Both span $y=0.05$--0.85 at $z=0.76$.\\
Devices and interval & 51 \texttt{THERMOCOUPLE} probes, 1~mm diameter, at $y=0.45$, $z=0.74$. Temperature rise averaged over 90--120~s. & \texttt{LAYER HEIGHT} on eight vertical columns spanning $z=0$--0.76 at $y=0.45$ (Section~2.2). CO and visibility at the same stations at $z=0.38$.\\
\end{reporttable}

\subhead{Boundary and burner settings shared by both runs}
The computational domain spans $x=0$--6.60, $y=0$--1.40 and $z=0$--0.76. The side wall occupies $y=0.90$--0.95. R1 is cut into it at $x=3.91$--4.34 and $z=0$--0.63. The exterior strip has \texttt{OPEN} boundaries at $y=1.40$ and on its ends and top over $y=0.95$--1.40. Both runs use the Deardorff gas-phase closure, end at 120~s and write device and HRR output at nominal 0.5~s intervals.

The floor burner spans $x=3.175$--3.425 and $y=0.325$--0.575 in the inputs. After alignment to the mesh it covers $5\times5$ horizontal cells of side 0.05~m in both runs, giving 0.0625~m$^2$. Both fire surfaces specify \texttt{HRRPUA} = 1699.2~kW/m$^2$ and use the FDS default startup factor $\tanh(t/1\,\mathrm{s})$~\cite{fdsuser}.

\subhead{Exhaust setting and recorded flow}
Each ceiling vent feeds an HVAC duct connected to an ambient node. Each duct has area 0.32~m$^2$, length 0.40~m and \texttt{VOLUME\_FLOW} = 0.56~m$^3$/s. The mean total recorded by the two \texttt{DUCT VOLUME FLOW} devices in this run was 1.11345~m$^3$/s over 0--120~s. Table~C3 reports the height check at the 1.12~m$^3$/s setpoint.

\clearpage
\reportheading{Appendix G.\quad PINN training and minimum-height correction}
The corrected PINN used in Section~5 fits the smoke-layer histories at eight FDS stations and then corrects the predicted minimum at each station. The model takes time and seven case inputs. All dimensions and flow rates in Table~G1 use m and m$^3$/s, respectively.

\tabletitle{Table G1. Case inputs to the Section~5 surrogate.}
\begin{reporttable}{L{.45}L{.55}}
\rowcolor{tablegrey}Input & Value supplied to the model\\
Left and right end-inlet heights, $a_L,a_R$ & $(a_L-0.44)/0.08$ and $(a_R-0.44)/0.08$\\
Door state & 0 for OP-1 and 1 for OP-6\\
Soot yield, $Y_s$ & $(Y_s-0.05)/0.05$\\
Curtain depth, $d$ & $(d-0.24)/0.08$\\
Total extract flow, $Q$ & $(Q-1.28)/0.32$\\
Vertical mesh & 0 for 19 cells and 1 for 38 cells\\
\end{reporttable}

\subhead{Time-series fit and upper-layer balance}
A constant term and 96 fixed tanh terms form the time basis. Each tanh term takes $(t-c)/s+\boldsymbol b^\mathsf{T}\boldsymbol u$, where $\boldsymbol u$ is the input vector. The centres $c$ are uniform on 0--120~s and the scales $s$ are log-uniform on 1--25~s. The input weights $\boldsymbol b$ are drawn from a zero-mean normal distribution with standard deviation 0.35. Two separate NumPy random generators, both seeded at 3322, supply the time parameters and input weights, respectively. Multiplying the basis by $1-\exp[-t/(2\,\mathrm{s})]$ gives an initial layer height of $H=0.76$~m. The output coefficients are fitted to FDS layer depths $H-h_j$ with ridge penalty 0.1.

Let $\bar d(t)$ be the spatial mean of the predicted layer depths, weighted by the length of each station’s interval, bounded by the midpoints between adjacent stations and by the end walls. The upper-layer balance used in the fit is
\[
 r(t)=\frac{\mathrm d\bar d}{\mathrm dt}
       +\frac{Q}{(6.60\,\mathrm m)(0.90\,\mathrm m)H}\,\bar d(t),
 \qquad r(t)\approx \overline{r}_{t\geq10\,\mathrm{s}} .
\]
Subtracting each run's mean residual over $t\geq10$~s eliminates the unknown constant net source. The fit penalises the resulting squared residual with weight 25 along the spatial-average direction, using unit Euclidean norm for that direction. Station-to-station differences are fitted from the FDS data.

\subhead{Correction at each station}
The raw station minimum is taken from the clipped height curve on a 0.5~s grid over 0--120~s. For each training run, the correction target is the observed FDS station minimum minus this fitted minimum. A Gaussian kernel $k(\boldsymbol u,\boldsymbol v)=\exp(-\|\boldsymbol u-\boldsymbol v\|^2/2)$ fits the centred correction targets, with ridge penalty 0.25 and the training mean added back. The predicted offset is added at each station and the result is clipped to $[0,H]$. The lowest corrected station value is the PINN prediction for the scenario. Each holdout excludes the test design and all its mesh variants from both the time-series fit and the correction. The direct regression uses the same kernel and penalty on the scenario minima. The hybrid averages the two predictions.
